\documentclass[runningheads]{llncs}

\usepackage[T1]{fontenc}
\usepackage{graphicx}
\usepackage{booktabs}
\usepackage{multirow}
\usepackage{subcaption}
\usepackage{amsmath,amssymb}
\usepackage{textcomp}
\usepackage{url}
\usepackage{makecell}
\usepackage{array}
\usepackage{siunitx}
\usepackage{float}
\usepackage{placeins}

\graphicspath{{figures/}}

\sisetup{
  round-mode=places,
  round-precision=2,
  detect-all
}

\begin{document}

\title{Contribution Title}

\author{First Author\inst{1}\orcidID{0000-1111-2222-3333} \and
Second Author\inst{2,3}\orcidID{1111-2222-3333-4444} \and
Third Author\inst{3}\orcidID{2222--3333-4444-5555}}

\authorrunning{F. Author et al.}

\institute{
Princeton University, Princeton NJ 08544, USA \and
Springer Heidelberg, Tiergartenstr. 17, 69121 Heidelberg, Germany
\email{lncs@springer.com}\\\\
\url{http://www.springer.com/gp/computer-science/lncs} \and
ABC Institute, Rupert-Karls-University Heidelberg, Heidelberg, Germany\\\\
\email{\{abc,lncs\}@uni-heidelberg.de}
}

\maketitle

\begin{abstract}
Selecting lightweight convolutional neural networks for edge vision based only on clean accuracy can lead to suboptimal deployment decisions, since practical performance also depends on robustness to image corruptions, latency, throughput, and model size. This paper presents a comparative evaluation of MCUNet, MobileNetV3-Small, and EfficientNet-B0 across three benchmarks with distinct characteristics: CIFAR-10, CIFAR-100, and Wake Vision person detection. We assess robustness using corruption benchmarks derived from the CIFAR-C protocol, comprising 15 corruption types at five severity levels, and we measure deployment efficiency with float32 TensorFlow Lite CPU-only inference on a Samsung Galaxy A14, a Samsung Galaxy S24+, and a Raspberry Pi Zero 2 W. The analysis reports clean accuracy, mean Corruption Error (mCE), relative mCE, P90 latency, throughput, model footprint, memory usage, and Friedman/Nemenyi statistical comparisons. Across the CIFAR benchmarks, EfficientNet-B0 achieves the highest clean accuracy and the lowest absolute mCE, whereas MobileNetV3-Small consistently provides the lowest latency and the smallest deployment footprint on all evaluated devices. On Wake Vision, MobileNetV3-Small also yields the best relative robustness, while MCUNet does not surpass the other models in the final float32 TensorFlow Lite deployment setting. Overall, the results indicate that edge-model selection should be guided by joint robustness--efficiency evidence rather than clean accuracy alone.
\keywords{Edge vision \and Lightweight CNNs \and Robustness \and TensorFlow Lite \and TinyML}
\end{abstract}

\section{Introduction}

Edge vision applications are increasingly deployed on resource-constrained platforms such as smartphones, single-board computers, and embedded devices. In these settings, model selection cannot be based on clean accuracy alone, since practical deployment also depends on robustness to input degradation, inference latency, throughput, and storage footprint. A model that performs well in offline evaluation may still be unsuitable in practice if it degrades under common corruptions or fails to meet runtime constraints on target hardware.

This issue is particularly relevant because lightweight convolutional neural networks are often compared under limited evaluation protocols. In many cases, models are selected from clean-test accuracy on a single dataset and then validated on only one deployment platform. Such a protocol provides an incomplete view of real-world suitability, especially for edge-vision scenarios in which images may be affected by noise, blur, compression artifacts, contrast shifts, or other non-adversarial perturbations.

A more reliable assessment should jointly consider predictive quality and deployment behavior. From the robustness perspective, common-corruption benchmarks provide a controlled way to evaluate sensitivity to realistic distribution shift. From the systems perspective, latency, throughput, and model footprint determine whether a network can be executed under practical edge constraints. Moreover, conclusions drawn from a single task may not generalize across different edge-vision regimes, such as multi-class natural-image classification and binary TinyML person detection.

In this paper, we conduct a comparative evaluation of three representative lightweight architectures with different design goals: MobileNetV3-Small, EfficientNet-B0, and MCUNet. We study them across three complementary benchmarks, namely CIFAR-10, CIFAR-100, and Wake Vision, and evaluate their behavior under a unified float32 TensorFlow Lite CPU-only deployment setting. Robustness is measured using corruption benchmarks derived from the CIFAR-C protocol, whereas efficiency is assessed through on-device inference on two Android smartphones and a Raspberry Pi Zero 2 W.

The main contributions of this paper are as follows:
\begin{enumerate}
    \item We provide a unified cross-benchmark comparison of three lightweight CNN architectures over CIFAR-10, CIFAR-100, and Wake Vision, combining clean accuracy, corruption robustness, and on-device performance.
    \item We report deployment-oriented efficiency metrics under a common float32 TensorFlow Lite CPU-only setup, including P90 latency, throughput, memory usage, and model footprint across heterogeneous edge devices.
    \item We analyze whether the observed differences are consistent by applying Friedman and Nemenyi statistical tests to robustness and device-level performance measurements.
\end{enumerate}

Our results show a clear trade-off between robustness and efficiency. EfficientNet-B0 achieves the strongest clean accuracy and absolute corruption robustness, while MobileNetV3-Small consistently offers the best runtime efficiency and the smallest deployment footprint. These findings suggest that edge-model selection should be treated as a multi-objective decision problem rather than as a ranking based solely on clean accuracy.

The remainder of this paper is organized as follows. Section~\ref{sec:related} reviews the related literature. Section~\ref{sec:method} describes the methodology and experimental setup. Section~\ref{sec:results} presents the empirical results. Section~\ref{sec:conclusion} concludes the paper and outlines directions for future work.

\section{Related Work}
\label{sec:related}
Lightweight CNNs for edge inference have been shaped by two major lines of work. The first focuses on mobile-friendly operator design, as in MobileNetV3, which combines depthwise separable convolutions, squeeze-and-excitation blocks, and neural architecture search to reduce compute cost while retaining accuracy~\cite{howard2019mobilenetv3}. The second focuses on compound scaling, as in EfficientNet, which jointly scales width, depth, and input resolution to improve accuracy-efficiency trade-offs~\cite{tan2019efficientnet}. For even tighter memory budgets, MCUNet targets microcontroller-class hardware through joint architecture and runtime co-design~\cite{lin2020mcunet}. Although these families are frequently discussed as ``efficient'', their relative behavior under corruption and heterogeneous deployment conditions is less often validated systematically.

Robustness to common corruptions has become a standard complement to clean accuracy. Hendrycks and Dietterich introduced the CIFAR-10-C and CIFAR-100-C benchmarks together with corruption-error metrics that quantify sensitivity to non-adversarial distribution shift~\cite{hendrycks2019benchmarking}. Later analyses emphasized that robustness is multi-faceted and should not be reduced to a single notion of generalization~\cite{hendrycks2021manyfaces}. RobustBench further contributed to the standardization of robustness evaluation workflows~\cite{croce2021robustbench}. At the representation level, Geirhos et al.~\cite{geirhos2019texture} argued that ImageNet-pretrained CNNs exhibit texture-biased behavior, while newer evidence suggests that the phenomenon depends strongly on task design and suppression protocol~\cite{burgert2025texture}.

Wake Vision extends this discussion to TinyML person detection by providing a large-scale benchmark tailored to edge computer vision~\cite{banbury2024wakevision}. Unlike CIFAR-10 and CIFAR-100, Wake Vision emphasizes binary person detection under realistic demographic and environmental diversity. It is therefore a useful complement to standard image-classification benchmarks when assessing whether architectural conclusions transfer across domains.

Finally, methodology matters. Fine-tuning choices can significantly alter the robustness inherited from pretraining~\cite{hwang2024pretraining}, while deployment performance is heavily shaped by runtime implementations such as TensorFlow Lite and XNNPACK~\cite{tensorflowlite,dukhan2021xnnpack}. Yet few studies jointly compare lightweight architectures under common corruptions, device-level measurements, and formal statistical tests over multiple datasets. This is the gap our work addresses.

\section{Methodology}
\label{sec:method}
\subsection{Experimental Pipeline}
Our experimental workflow comprised four stages: (i) supervised training and
checkpoint selection, (ii) corruption-based robustness evaluation and
aggregation of robustness metrics, (iii) TensorFlow Lite export, and (iv)
on-device benchmarking on Android smartphones and Raspberry Pi. Training and
analysis were conducted with TensorFlow/Keras notebooks, smartphone
benchmarking was implemented in a Kotlin/Jetpack Compose Android application,
and Raspberry Pi evaluation was executed with a Python-based inference script.
For double-blind review, the complete reproducibility package is available in
an anonymous repository.

\subsection{Datasets and Corruption Protocol}
We selected three datasets that represent complementary edge-vision scenarios.
\begin{itemize}
    \item \textbf{CIFAR-10}~\cite{krizhevsky2009cifar}: 10 natural-image
    classes. Inputs were resized to $160\times160$, and the last 5,000
    training images were reserved for validation.

    \item \textbf{CIFAR-100}~\cite{krizhevsky2009cifar}: 100 classes with
    higher inter-class difficulty than CIFAR-10. The same $160\times160$
    preprocessing and validation split were adopted.

    \item \textbf{Wake Vision}~\cite{banbury2024wakevision}: binary
    person-versus-no-person classification. We used the official
    \texttt{train\_quality}, \texttt{validation}, and \texttt{test} splits
    provided through \texttt{tensorflow\_datasets}, resized inputs to
    $224\times224$, and excluded ambiguous \texttt{person == -1} samples from
    the test split during standard evaluation.
\end{itemize}

To assess robustness, we followed the common-corruption evaluation rationale
of Hendrycks and Dietterich and considered the 15 corruption types used in the
CIFAR-C family, each evaluated at five severity
levels~\cite{hendrycks2019benchmarking}. This results in 75 corrupted
conditions per dataset. For CIFAR-10 and CIFAR-100, we reused the corruption
outputs already available in the repository. For Wake Vision, we generated a
matched corruption benchmark by applying the same 15 common corruption types at
five severity levels to the clean test split, using the
\texttt{imagecorruptions} library~\cite{michaelis2019imagecorruptions}. Unlike
the CIFAR-C datasets, which are distributed as pre-generated \texttt{.npy}
arrays, corruptions were applied on-the-fly per image and immediately
discarded, since pre-generating all 75 corrupted conditions for approximately
55{,}762 images at $224\times224$ would require approximately 780\,GB of
storage.

\subsection{Architectures and Training Strategy}

We compared three lightweight convolutional architectures with distinct design
goals: \textbf{MobileNetV3-Small}, a mobile-oriented backbone scaled for
efficiency; \textbf{EfficientNet-B0}, a compound-scaled backbone balancing
accuracy and parameter count; and \textbf{MCUNet}, originally designed for
microcontroller deployment, converted from PyTorch to Keras prior to
fine-tuning. All three were initialized from ImageNet-pretrained weights and
adapted with a task-specific classification head.

To ensure that observed differences reflect architectural properties rather
than pipeline variation, all models were trained under a strictly shared
protocol implemented in a single function
(\texttt{train\_head\_and\_finetune}). Table~\ref{tab:training_protocol}
summarizes the shared hyperparameters and the architecture-specific
preprocessing conventions.

\begin{table}[ht]
\centering
\caption{Shared training protocol and per-architecture preprocessing.}
\label{tab:training_protocol}
\begin{tabular}{ll}
\toprule
\textbf{Hyperparameter / Setting} & \textbf{Value} \\
\midrule
\multicolumn{2}{l}{\textit{Shared across all architectures}} \\
\midrule
Input resolution        & $160\times160$ (CIFAR-10/100);
                          $224\times224$ (Wake Vision) \\
Optimizer               & AdamW ($\lambda=10^{-4}$), Adam as fallback \\
Head learning rate      & $10^{-3}$ \\
Fine-tune learning rate & $10^{-5}$ \\
Label smoothing         & $\varepsilon = 0.1$ \\
Early stopping          & Patience 8, monitors val.\ accuracy \\
LR reduction on plateau & Patience 3, monitors val.\ accuracy \\
Data augmentation       & None (\texttt{USE\_AUGMENTATION = False}) \\
Fine-tune ratio         & Top 20\% of backbone layers frozen; \\
                        & BatchNorm layers kept in inference mode \\
Model selection         & Best validation accuracy
                          (\texttt{ModelCheckpoint}) \\
Export precision        & float32, no quantization \\
\midrule
\multicolumn{2}{l}{\textit{Per-architecture preprocessing}} \\
\midrule
MobileNetV3-Small & Keras rescaling to $[-1, 1]$ \\
EfficientNet-B0   & ImageNet-compatible preprocessing (Keras built-in) \\
MCUNet            & Explicit ImageNet normalization:
                    $(x/255 - \mu)/\sigma$ \\
\bottomrule
\end{tabular}
\end{table}

Training proceeded in two stages. In the first stage, only the classification
head was optimized with the backbone fully frozen. In the second stage, the
final 80\% of backbone layers were unfrozen whenever the architecture exposed
a separable backbone, while all \texttt{BatchNormalization} layers remained in
inference mode throughout. The complete training code is available in the
anonymous repository.

Each model was trained with a single fixed random seed (\texttt{seed\,=\,42}),
set via \texttt{tf.keras.utils.set\_random\_seed} prior to data loading and
model construction. No repeated training runs were performed; consequently,
the reported accuracy and robustness figures reflect a single-seed outcome per
architecture and dataset. Variance across seeds is therefore not estimated in
this study. The Friedman and Nemenyi statistical tests described in
Section~\ref{sec:method} are applied over the 45 dataset--corruption blocks
(three datasets $\times$ 15 corruption types), not over training repetitions;
their purpose is to assess whether the observed \emph{ranking} of models is
consistent across corruption conditions rather than across independent training
runs.

\subsection{Robustness Metrics}
Let $Acc_{m,c,s}$ denote the accuracy of model $m$ under corruption type $c$
at severity level $s$, and let $E_{m,c,s}=1-Acc_{m,c,s}$ denote the
corresponding error. Following the corruption-error framework introduced by
Hendrycks and Dietterich, corruption robustness is evaluated by aggregating
errors across severity levels and normalizing them with respect to a reference
model under the same corruption~\cite{hendrycks2019benchmarking}.

The corruption error of model $m$ relative to baseline model $b$ for
corruption $c$ is defined as:
\begin{equation}
CE_{m,c}^{(b)} =
\frac{\sum_{s=1}^{5} E_{m,c,s}}
     {\sum_{s=1}^{5} E_{b,c,s}}.
\end{equation}

The mean Corruption Error (mCE) is the arithmetic mean over the full
corruption set $C$:
\begin{equation}
mCE_{m}^{(b)} = \frac{1}{|C|}\sum_{c \in C} CE_{m,c}^{(b)}.
\end{equation}

In addition to mCE, we report a relative corruption measure that discounts
the contribution of clean error, following the motivation of Relative CE in
the original benchmark~\cite{hendrycks2019benchmarking}. Let
$E_{m,\mathrm{clean}}$ denote the clean error of model $m$. We define:
\begin{equation}
RCE_{m,c}^{(b)} =
\frac{\sum_{s=1}^{5} E_{m,c,s} - E_{m,\mathrm{clean}}}
     {\sum_{s=1}^{5} E_{b,c,s} - E_{b,\mathrm{clean}}}
\end{equation}
and compute relative mCE as
\begin{equation}
RmCE_{m}^{(b)} = \frac{1}{|C|}\sum_{c \in C} RCE_{m,c}^{(b)}.
\end{equation}

Unlike the canonical ImageNet-C protocol, which reports corruption error
relative to a fixed external reference such as AlexNet, our study adopts an
internal baseline among the evaluated
architectures~\cite{hendrycks2019benchmarking}. We use MCUNet as the reference
model because it consistently yielded the highest clean error across the
considered datasets in our experiments. This choice provides a conservative
and interpretable normalization within the comparative setting: values below
100 indicate better corruption performance than the internal baseline, whereas
MCUNet is fixed at 100 by definition.

\subsection{On-device Benchmarking}

Deployment efficiency was evaluated under a unified CPU-only TensorFlow Lite
setting across all devices, with no hardware delegate enabled and no
post-training quantization applied. Beyond latency and throughput, both
benchmarks report memory usage, enabling a direct resource-efficiency
comparison across the three evaluated platforms.

On Android, we implemented a custom benchmark application in Kotlin and Jetpack
Compose. For each dataset, the application loads the corresponding TFLite
model and assets, performs 10 warm-up inferences, and then measures per-image
inference time via \texttt{Interpreter.run()} with up to 4 CPU threads. The
benchmark reports mean latency, P90 latency, standard deviation, throughput,
accuracy, and memory usage over a balanced subset of 400 images per dataset.
We emphasize P90 latency throughout the analysis because it captures tail
behavior more faithfully than the mean, which is more informative for
latency-sensitive edge applications.

For the Raspberry Pi Zero 2 W, we used a Python benchmark based on
\texttt{tflite\_runtime.interpreter} (with fallback to
\texttt{tf.lite.Interpreter}), executed inside Docker to control the execution
environment. All runs were performed with \texttt{num\_threads=1}. Although
the device exposes four Cortex-A53 cores, multi-threaded TFLite inference on
low-power ARM boards incurs synchronization overhead that can negate the
expected parallelism gain~\cite{tflite_best_practices}. More importantly,
restricting inference to a single thread places the Pi Zero 2 W in the same
operating regime as single-core embedded processors (e.g., Cortex-M4/M7
class), making it a conservative proxy for the most resource-constrained tier
of edge deployment. The same 400-image balanced subsets and 10 warm-up runs
used on Android were preserved for comparability.

\subsection{Statistical Testing}
To assess whether the observed differences were consistent across evaluation
conditions, we applied Friedman and Nemenyi analyses following Dem\v{s}ar's
recommendations for comparing multiple models over repeated
blocks~\cite{demsar2006statistical}. For robustness, each dataset-corruption
pair defined one block, and the models were ranked according to corruption
error, with lower values indicating better robustness. Across the three
datasets and 15 corruption types, this yields 45 robustness blocks.

For device efficiency, each dataset-device pair defined one block. Models were
ranked separately according to P90 latency, where lower is better, and
throughput, where higher is better. Considering three datasets and three
devices, this yields 9 device-level blocks. This statistical design allows us
to test whether the ordering observed in the descriptive results remains stable
across corruption types and heterogeneous deployment settings.

\section{Results}
\label{sec:results}

\subsection{Model Size}
Table~\ref{tab:model_size} reports the storage footprint of the exported models. Model size was primarily determined by architecture and varied only slightly across datasets. MobileNetV3-Small was consistently the most compact model, ranging from 3.58 to 3.78 MB in TensorFlow Lite format, MCUNet occupied an intermediate range from 5.39 to 5.50 MB, and EfficientNet-B0 was the largest model, ranging from 15.33 to 15.79 MB.

These differences are substantial from a deployment perspective. In TensorFlow Lite format, EfficientNet-B0 required approximately $4.2\times$ the storage of MobileNetV3-Small and approximately $2.8\times$ that of MCUNet. Therefore, even before runtime measurements are considered, EfficientNet-B0 imposes a markedly higher storage cost on memory-constrained edge targets.

\begin{table}[H]
\caption{Model storage footprint across datasets.}
\label{tab:model_size}
\centering
\scriptsize
\setlength{\tabcolsep}{3pt}
\resizebox{\textwidth}{!}{%
\begin{tabular}{llrrrr}
\toprule
Dataset & Model & Keras (MB) & Keras.zip (MB) & TFLite (MB) & TFLite.zip (MB) \\
\midrule
CIFAR-10 & MobileNetV3-Small & 4.14 & 3.40 & 3.58 & 3.31 \\
CIFAR-10 & EfficientNet-B0 & 16.30 & 14.42 & 15.35 & 14.21 \\
CIFAR-10 & MCUNet & 6.02 & 5.15 & 5.39 & 5.01 \\
CIFAR-100 & MobileNetV3-Small & 4.34 & 3.58 & 3.78 & 3.49 \\
CIFAR-100 & EfficientNet-B0 & 16.74 & 14.83 & 15.79 & 14.62 \\
CIFAR-100 & MCUNet & 6.13 & 5.25 & 5.50 & 5.11 \\
Wake Vision & MobileNetV3-Small & 4.13 & 3.38 & 3.59 & 3.29 \\
Wake Vision & EfficientNet-B0 & 16.26 & 14.38 & 15.33 & 14.17 \\
Wake Vision & MCUNet & 6.01 & 5.14 & 5.40 & 5.00 \\
\bottomrule
\end{tabular}%
}
\end{table}

\subsection{On-device Efficiency}
Figures~\ref{fig:mobile_latency} and~\ref{fig:mobile_throughput} summarize the smartphone results. The ordering was consistent across both phones and all three datasets: MobileNetV3-Small was always the fastest model, MCUNet was consistently second, and EfficientNet-B0 was always the slowest. This pattern indicates that the efficiency ranking was stable not only across tasks but also across mobile devices with distinct performance tiers.

On the Galaxy S24+, MobileNetV3-Small achieved P90 latency between 1.83 and 4.07 ms, whereas EfficientNet-B0 ranged from 12.87 to 36.11 ms. On the Galaxy A14, the same ranking was preserved, but the absolute gap became larger: MobileNetV3-Small remained below 9 ms across all datasets, MCUNet ranged from 16.87 to 33.40 ms, and EfficientNet-B0 reached 77.15 ms on Wake Vision. Throughput followed the inverse trend, with MobileNetV3-Small peaking at 158.63 images/s on CIFAR-10 on the Galaxy S24+ and still sustaining 30.91 images/s on Wake Vision on the Galaxy A14.

\begin{figure}[H]
\centering
\includegraphics[width=\textwidth]{mobile_p90_latency.pdf}
\caption{P90 latency on Android phones. Lower is better.}
\label{fig:mobile_latency}
\end{figure}

\begin{figure}[H]
\centering
\includegraphics[width=\textwidth]{mobile_throughput.pdf}
\caption{Throughput on Android phones. Higher is better.}
\label{fig:mobile_throughput}
\end{figure}

When the Raspberry Pi Zero 2 W is included, the device effect becomes more pronounced, as shown in Figures~\ref{fig:edge_latency} and~\ref{fig:edge_throughput}. MobileNetV3-Small remained the only model with clearly interactive latency across all evaluated tasks, ranging from 25.27 to 51.29 ms. MCUNet retained the second-best position, but its latency increased sharply on the single-board computer, reaching 182.42 ms on CIFAR-100 and 465.81 ms on Wake Vision. EfficientNet-B0 incurred the largest runtime penalty, exceeding 290 ms on both CIFAR variants and reaching 544.01 ms on Wake Vision.

The Raspberry Pi benchmark was executed with a single thread, one model at a time for each dataset, and inside Docker, which makes the results more directly comparable with the Android measurements while controlling for runtime variability. The same pattern can be observed in throughput. On the Raspberry Pi, MobileNetV3-Small sustained between 8.72 and 34.67 images/s depending on the dataset, whereas EfficientNet-B0 dropped to 1.68--3.45 images/s. Memory usage on the Raspberry Pi was also captured in the updated benchmark, and it followed the same general ordering as runtime: EfficientNet-B0 used the most memory, MobileNetV3-Small the least, and MCUNet remained intermediate in most settings.

\begin{figure}[H]
\centering
\includegraphics[width=\textwidth]{edge_p90_latency.pdf}
\caption{P90 latency across all edge devices, including the Raspberry Pi Zero 2 W. Lower is better.}
\label{fig:edge_latency}
\end{figure}

\begin{figure}[H]
\centering
\includegraphics[width=\textwidth]{edge_throughput.pdf}
\caption{Throughput across all edge devices, including the Raspberry Pi Zero 2 W. Higher is better.}
\label{fig:edge_throughput}
\end{figure}

Table~\ref{tab:device_efficiency} consolidates the detailed device-level measurements. Two additional observations are relevant. First, on-device accuracies remained close to the clean-set ranking observed offline, which suggests that float32 TensorFlow Lite export preserved the relative predictive behavior of the models. Second, Android memory consumption generally increased with model size, with EfficientNet-B0 usually requiring the highest memory footprint and MobileNetV3-Small the lowest. On the Raspberry Pi, memory usage was now measured explicitly and was consistently higher for EfficientNet-B0 than for the other models. The Raspberry Pi benchmark was executed in Docker and with a single thread, which improves fairness but also makes the absolute values more representative of a controlled CPU-only deployment than of a fully optimized native setup.

\begin{table}[p]
\caption{Detailed on-device efficiency results.}
\label{tab:device_efficiency}
\centering
\scriptsize
\setlength{\tabcolsep}{2pt}
\resizebox{\textwidth}{!}{%
\begin{tabular}{lllrrrr}
\toprule
Device & Dataset & Model & P90 (ms) & Throughput & Acc. (\%) & Mem. (MB) \\
\midrule
S24+ & CIFAR-10 & MobileNetV3-Small & 1.83 & 158.63 & 89.25 & 22.21 \\
S24+ & CIFAR-10 & MCUNet & 5.98 & 102.37 & 73.50 & 41.89 \\
S24+ & CIFAR-10 & EfficientNet-B0 & 12.87 & 59.94 & 94.00 & 75.72 \\
S24+ & CIFAR-100 & MobileNetV3-Small & 2.05 & 152.26 & 75.75 & 35.11 \\
S24+ & CIFAR-100 & MCUNet & 7.93 & 85.44 & 53.25 & 49.19 \\
S24+ & CIFAR-100 & EfficientNet-B0 & 17.04 & 47.08 & 82.00 & 61.29 \\
S24+ & Wake Vision & MobileNetV3-Small & 4.07 & 53.52 & 89.50 & 22.03 \\
S24+ & Wake Vision & MCUNet & 17.18 & 32.81 & 84.00 & 37.59 \\
S24+ & Wake Vision & EfficientNet-B0 & 36.11 & 20.80 & 90.50 & 70.15 \\
A14 & CIFAR-10 & MobileNetV3-Small & 4.53 & 69.95 & 89.25 & 19.89 \\
A14 & CIFAR-10 & MCUNet & 16.87 & 38.54 & 73.50 & 28.42 \\
A14 & CIFAR-10 & EfficientNet-B0 & 39.85 & 21.45 & 94.00 & 45.59 \\
A14 & CIFAR-100 & MobileNetV3-Small & 4.58 & 71.20 & 75.75 & 22.65 \\
A14 & CIFAR-100 & MCUNet & 16.91 & 38.69 & 53.25 & 25.49 \\
A14 & CIFAR-100 & EfficientNet-B0 & 39.60 & 20.93 & 82.00 & 50.13 \\
A14 & Wake Vision & MobileNetV3-Small & 8.85 & 30.91 & 89.50 & 26.83 \\
A14 & Wake Vision & MCUNet & 33.40 & 18.26 & 84.00 & 19.67 \\
A14 & Wake Vision & EfficientNet-B0 & 77.15 & 10.73 & 90.50 & 43.79 \\
RPi Zero 2 W & CIFAR-10 & MobileNetV3-Small & 25.63 & 34.06 & 87.25 & 15.80 \\
RPi Zero 2 W & CIFAR-10 & MCUNet & 182.70 & 5.39 & 86.25 & 21.10 \\
RPi Zero 2 W & CIFAR-10 & EfficientNet-B0 & 350.34 & 2.83 & 94.50 & 20.65 \\
RPi Zero 2 W & CIFAR-100 & MobileNetV3-Small & 25.75 & 33.90 & 72.00 & 16.10 \\
RPi Zero 2 W & CIFAR-100 & MCUNet & 182.96 & 5.38 & 51.75 & 21.20 \\
RPi Zero 2 W & CIFAR-100 & EfficientNet-B0 & 350.85 & 2.83 & 81.75 & 20.75 \\
RPi Zero 2 W & Wake Vision & MobileNetV3-Small & 51.02 & 8.69 & 89.50 & 18.06 \\
RPi Zero 2 W & Wake Vision & MCUNet & 462.20 & 1.91 & 85.25 & 23.21 \\
RPi Zero 2 W & Wake Vision & EfficientNet-B0 & 663.06 & 1.39 & 91.00 & 37.44 \\
\bottomrule
\end{tabular}%
}
\end{table}

The statistical analysis supports the same ordering. Over the nine dataset-device blocks, Friedman tests yielded $\chi^2_F = 18.00$ with $p = 1.23\times10^{-4}$ for P90 latency, and the same result was obtained for throughput. Because both metrics induce the same ranking, Figure~\ref{fig:cd_edge_latency} reports only the Nemenyi critical-difference diagram for P90 latency, which ranks MobileNetV3-Small first, MCUNet second, and EfficientNet-B0 third. A finer-grained run-level latency analysis over 300 blocks led to the same ranking with substantially stronger evidence ($\chi^2_F = 453.84$, $p \approx 2.8\times10^{-99}$), indicating that the efficiency ordering is highly stable under the evaluated deployment protocol.

\begin{figure}[H]
\centering
\includegraphics[width=0.72\textwidth]{cd_edge_p90_latency.pdf}
\caption{Nemenyi critical-difference diagram for edge-device P90 latency.}
\label{fig:cd_edge_latency}
\end{figure}

\subsection{Robustness}
Figures~\ref{fig:robustness_c10}, \ref{fig:robustness_c100}, and~\ref{fig:robustness_wake} summarize the trade-off between clean accuracy and corruption robustness. Across all three datasets, EfficientNet-B0 achieved both the highest clean accuracy and the lowest absolute mCE, while MobileNetV3-Small consistently ranked second and MCUNet ranked third. The gap was modest on CIFAR-10-C, but became more pronounced on CIFAR-100-C and Wake Vision-C.

This trend is reflected in Table~\ref{tab:robustness_summary}. On CIFAR-10, EfficientNet-B0 reduced mCE to 42.82, compared with 58.95 for MobileNetV3-Small and 100.00 for MCUNet. On CIFAR-100, the same ordering was preserved, with mCE values of 66.33, 76.32, and 100.00, respectively. On Wake Vision, EfficientNet-B0 again achieved the best absolute corruption performance with mCE equal to 58.72, while MobileNetV3-Small reached 71.91.

\begin{figure}[H]
\centering
\includegraphics[width=\textwidth]{cifar10c_architecture_corruption_robustness.pdf}
\caption{Clean-accuracy and corruption-robustness comparison on CIFAR-10-C.}
\label{fig:robustness_c10}
\end{figure}

\begin{figure}[H]
\centering
\includegraphics[width=\textwidth]{cifar100c_architecture_corruption_robustness.pdf}
\caption{Clean-accuracy and corruption-robustness comparison on CIFAR-100-C.}
\label{fig:robustness_c100}
\end{figure}

\begin{figure}[H]
\centering
\includegraphics[width=\textwidth]{wakevision_architecture_corruption_robustness.pdf}
\caption{Clean-accuracy and corruption-robustness comparison on Wake Vision-C.}
\label{fig:robustness_wake}
\end{figure}

Table~\ref{tab:robustness_summary} reports clean error, mCE, and relative mCE. Absolute mCE favors EfficientNet-B0 on all three datasets. Relative mCE preserves the same ordering on CIFAR-10 and CIFAR-100, but differs on Wake Vision, where MobileNetV3-Small attains a lower relative mCE than EfficientNet-B0. This indicates that, although EfficientNet-B0 remains superior in absolute corruption error, MobileNetV3-Small experiences a smaller corruption-related degradation relative to its own clean error on that task.

\begin{table}[t]
\caption{Clean error, mCE, and relative mCE (\%) with MCUNet as the baseline. Lower is better for all metrics.}
\label{tab:robustness_summary}
\centering
\small
\setlength{\tabcolsep}{4pt}
\begin{tabular}{llrrr}
\toprule
Dataset & Model & Clean Error & mCE & Rel.\ mCE \\
\midrule
CIFAR-10 & MobileNetV3-Small & 9.56 & 58.95 & 74.58 \\
CIFAR-10 & EfficientNet-B0 & 5.39 & 42.82 & 62.52 \\
CIFAR-10 & MCUNet & 21.12 & 100.00 & 100.00 \\
CIFAR-100 & MobileNetV3-Small & 26.99 & 76.32 & 98.36 \\
CIFAR-100 & EfficientNet-B0 & 21.68 & 66.33 & 94.75 \\
CIFAR-100 & MCUNet & 41.32 & 100.00 & 100.00 \\
Wake Vision & MobileNetV3-Small & 13.87 & 71.91 & 35.01 \\
Wake Vision & EfficientNet-B0 & 12.27 & 58.72 & 63.11 \\
Wake Vision & MCUNet & 18.73 & 100.00 & 100.00 \\
\bottomrule
\end{tabular}
\end{table}

\begin{table}[p]
\caption{Per-corruption CE (\%) on CIFAR-10-C. MCUNet is the baseline and therefore fixed at 100. Lower is better.}
\label{tab:cifar10_ce}
\centering
\scriptsize
\begin{tabular}{lrrr}
\toprule
Corruption & MobileNetV3-Small & EfficientNet-B0 & MCUNet \\
\midrule
Gaussian Noise & 72 & 53 & 100 \\
Shot Noise & 68 & 49 & 100 \\
Impulse Noise & 90 & 46 & 100 \\
Defocus Blur & 49 & 34 & 100 \\
Glass Blur & 74 & 82 & 100 \\
Motion Blur & 58 & 45 & 100 \\
Zoom Blur & 49 & 36 & 100 \\
Snow & 45 & 37 & 100 \\
Frost & 39 & 29 & 100 \\
Fog & 78 & 47 & 100 \\
Brightness & 33 & 19 & 100 \\
Contrast & 75 & 55 & 100 \\
Elastic Transform & 57 & 47 & 100 \\
Pixelate & 54 & 39 & 100 \\
JPEG Compression & 53 & 43 & 100 \\
\bottomrule
\end{tabular}
\end{table}

\begin{table}[p]
\caption{Per-corruption CE (\%) on CIFAR-100-C. MCUNet is the baseline and therefore fixed at 100. Lower is better.}
\label{tab:cifar100_ce}
\centering
\scriptsize
\begin{tabular}{lrrr}
\toprule
Corruption & MobileNetV3-Small & EfficientNet-B0 & MCUNet \\
\midrule
Gaussian Noise & 87 & 76 & 100 \\
Shot Noise & 84 & 73 & 100 \\
Impulse Noise & 101 & 69 & 100 \\
Defocus Blur & 69 & 60 & 100 \\
Glass Blur & 89 & 99 & 100 \\
Motion Blur & 74 & 68 & 100 \\
Zoom Blur & 70 & 62 & 100 \\
Snow & 68 & 62 & 100 \\
Frost & 59 & 55 & 100 \\
Fog & 85 & 69 & 100 \\
Brightness & 55 & 44 & 100 \\
Contrast & 81 & 73 & 100 \\
Elastic Transform & 76 & 70 & 100 \\
Pixelate & 70 & 62 & 100 \\
JPEG Compression & 76 & 67 & 100 \\
\bottomrule
\end{tabular}
\end{table}

\begin{table}[p]
\caption{Per-corruption CE (\%) on Wake Vision-C. MCUNet is the baseline and therefore fixed at 100. Lower is better.}
\label{tab:wake_ce}
\centering
\scriptsize
\begin{tabular}{lrrr}
\toprule
Corruption & MobileNetV3-Small & EfficientNet-B0 & MCUNet \\
\midrule
Gaussian Noise & 68 & 48 & 100 \\
Shot Noise & 72 & 51 & 100 \\
Impulse Noise & 70 & 47 & 100 \\
Defocus Blur & 66 & 72 & 100 \\
Glass Blur & 62 & 51 & 100 \\
Motion Blur & 69 & 61 & 100 \\
Zoom Blur & 78 & 73 & 100 \\
Snow & 63 & 48 & 100 \\
Frost & 68 & 58 & 100 \\
Fog & 90 & 59 & 100 \\
Brightness & 69 & 55 & 100 \\
Contrast & 100 & 70 & 100 \\
Elastic Transform & 71 & 67 & 100 \\
Pixelate & 65 & 60 & 100 \\
JPEG Compression & 68 & 61 & 100 \\
\bottomrule
\end{tabular}
\end{table}

The statistical test over the 45 dataset-corruption blocks confirms that these differences are not isolated cases. Friedman yielded $\chi^2_F = 104.46$ with $p = 2.08\times10^{-23}$, and the Nemenyi diagram in Figure~\ref{fig:robustness_cd} ranks EfficientNet-B0 first, MobileNetV3-Small second, and MCUNet third. In other words, the robustness ordering is as stable as the efficiency ordering, but it favors the largest model rather than the most lightweight one.

\begin{figure}[H]
\centering
\includegraphics[width=0.72\textwidth]{cd_robustness_ce.pdf}
\caption{Nemenyi critical-difference diagram for corruption error.}
\label{fig:robustness_cd}
\end{figure}

\subsection{Efficiency--Robustness Trade-off}
Taken together, the results indicate a clear trade-off between predictive robustness and deployment efficiency under the evaluated setting.

\begin{itemize}
    \item \textbf{MobileNetV3-Small offers the best deployment-oriented balance.} It is the smallest model, achieves the best latency and throughput on every evaluated device, and remains consistently second-best in absolute robustness across all three corruption benchmarks.
    \item \textbf{EfficientNet-B0 is the strongest model in predictive terms.} It delivers the best clean accuracy and the lowest absolute mCE on all datasets, but at the cost of substantially higher storage, latency, and memory usage.
    \item \textbf{MCUNet is less competitive in the evaluated deployment protocol.} It is smaller and faster than EfficientNet-B0, but it remains clearly behind MobileNetV3-Small in efficiency and behind both alternatives in robustness under the float32 TensorFlow Lite CPU-only setting considered here.
\end{itemize}

These findings should be interpreted within the scope of the present evaluation. MCUNet was originally designed for highly constrained TinyML and microcontroller scenarios, where memory scheduling, quantization, and specialized runtimes play a central role in performance~\cite{lin2020mcunet}. Therefore, its behavior in this work should be understood as specific to the chosen export and execution pipeline rather than as a universal statement about its suitability for all edge settings.

\section{Conclusion and Future Work}
\label{sec:conclusion}

This paper presented a cross-benchmark evaluation of lightweight CNNs for edge vision, combining clean accuracy, corruption robustness, model footprint, and on-device efficiency on two Android smartphones and a Raspberry Pi Zero 2 W. Across CIFAR-10, CIFAR-100, and Wake Vision, EfficientNet-B0 consistently achieved the strongest predictive performance, with the highest clean accuracy and the lowest absolute corruption error. In contrast, MobileNetV3-Small consistently provided the best deployment efficiency, combining the smallest storage footprint with the lowest latency, highest throughput, and lowest memory usage on the evaluated devices.

From a practical perspective, the results show that edge-model selection should be treated as a multi-objective decision problem. If predictive performance under common corruptions is the main priority, EfficientNet-B0 is the strongest option among the evaluated architectures. If the goal is a more balanced deployment profile under tight runtime, memory, and storage constraints, MobileNetV3-Small offers the most favorable trade-off in the studied float32 TensorFlow Lite CPU-only scenario. Under the same setting, MCUNet did not emerge as the preferred choice in either robustness or runtime efficiency.

These conclusions should be interpreted in light of the experimental scope. The study intentionally focused on a unified float32 TensorFlow Lite CPU-only deployment setting in order to enable a controlled comparison across heterogeneous devices. On the Raspberry Pi Zero 2 W, the benchmark was further restricted to a single thread, each model was executed independently for each dataset, and the runs were also repeated inside Docker to improve experimental fairness and environmental control. As a result, the findings do not necessarily reflect the behavior of the same architectures under quantized inference, hardware delegates, or microcontroller-specific runtimes, which are especially relevant for models such as MCUNet~\cite{lin2020mcunet}.

Future work should extend the analysis in three directions. First, the same comparison should be repeated under quantized inference and hardware-accelerated execution, including GPU, NPU, and vendor-specific delegates, to determine whether the observed ordering persists outside the CPU-only regime. Second, the architectural space should be expanded to include newer lightweight families and training strategies, such as corruption-aware augmentation, robust fine-tuning, pruning, and distillation. Third, the deployment evaluation should incorporate additional systems-level metrics, including energy consumption, thermal behavior, peak memory, and longer-running workloads, so that model selection can be grounded in a broader characterization of real edge operation.

\bibliographystyle{splncs04}
\bibliography{bib}

\end{document}